\section{Relational outputs and finite coherent selection}

Many constructive failures arise from asking for a single exact output where
the mathematics supplies only a finite family of locally valid outputs.
Relational presentations make this explicit.

\subsection{Finite candidate systems}

Fix a requested final depth \(K\).  For each level \(k\le K\), let
\(\Cand(k)\) be a finite coded list, and let
\(\Compat(k,c,d)\in\B\) mean that the level-\(k+1\) candidate \(d\) refines
the level-\(k\) candidate \(c\).  Membership and compatibility are
primitive-recursive, and all list lengths have primitive-recursive bounds.

There are two useful hypotheses.

\begin{description}[leftmargin=2.7em,style=nextline]
\item[Strong seriality.]
Every candidate in \(\Cand(k)\) has a compatible successor in
\(\Cand(k+1)\).
\item[Rooted viability.]
At least one candidate in \(\Cand(0)\) extends through level \(K\), but
dead-end candidates may also be present.
\end{description}

Strong seriality permits greedy forward selection.  Rooted viability is
strictly more general and is handled by backward dynamic programming.

\subsection{Backward viability}

Define a Boolean \(\mathsf{Viable}(K,k,c)\) by downward recursion:
\[
\begin{aligned}
 \mathsf{Viable}(K,K,c)
   &\defeq [c\in\Cand(K)],\\
 \mathsf{Viable}(K,k,c)
   &\defeq [c\in\Cand(k)]\wedge
      \anyfold_{d\in\Cand(k+1)}
      \bigl(\Compat(k,c,d)\wedge
            \mathsf{Viable}(K,k+1,d)\bigr).
\end{aligned}
\]
The recursion is implemented with fuel \(K-k\), so it is primitive recursive.
For a fixed \(K\), the finite table of viability bits is computed from the
last level backward.

\begin{theorem}[Finite-horizon coherent selection]
Suppose some \(c_0\in\Cand(0)\) is viable to level \(K\).  Then a primitive-
recursive bounded search produces
\[
    c_0,c_1,\ldots,c_K
\]
such that \(c_k\in\Cand(k)\) and
\(\Compat(k,c_k,c_{k+1})=1\) for every \(k<K\).
The correctness proof uses only quantifier-free induction.
\end{theorem}

\begin{proof}
Compute the viability table.  Choose the least viable element of
\(\Cand(0)\) by bounded search.  Having chosen a viable \(c_k\), the defining
equation for viability says that the finite existential fold at level
\(k+1\) is \(1\).  The bounded witness lemma returns the least \(d\) for
which both compatibility and viability hold; set \(c_{k+1}=d\).

The induction invariant at stage \(j\le K\) states:
the output list has length \(j+1\); its last member is viable to \(K\);
every member belongs to its candidate list; and all adjacent pairs are
compatible.  The invariant is a conjunction of Boolean folds over a list of
length at most \(K+1\), hence quantifier-free.  The step follows from the
witness property of viability.  At \(j=K\), viability reduces to membership.
\end{proof}

\begin{remark}[No hidden weak König lemma]
The theorem is finite in the free parameter \(K\).  It does not assert a
single infinite path through every infinite finitely branching tree.  For
each requested precision, it computes only the finite compatibility chain
needed at that precision.  An infinite regular name requires an additional
deterministic refinement operator, nested isolating regions, or a quantitative
uniqueness argument.
\end{remark}

\subsection{Compatibility as the semantic output}

Suppose a mathematical object \(x\) is represented at level \(k\) by a finite
family of regions \(R(k,c)\).  A candidate list is correct when the union of
those regions covers every object satisfying the input contract.  A
compatibility witness says that \(R(k+1,d)\subseteq R(k,c)\), or more
generally that the two observations describe the same object to the stated
tolerances.

This is often the right final theorem.  Examples include:
\begin{itemize}[leftmargin=1.6em]
\item finite disks covering all polynomial roots, with multiplicity data;
\item the two unit quaternions \(q,-q\) representing one rotation;
\item a projective line rather than a signed unit eigenvector;
\item a finite set of local analytic charts;
\item several nearly minimizing grid cells rather than an exact minimizer.
\end{itemize}

The framework never discards the relation and silently picks a member.  A
later consumer either accepts relational input, or supplies enough evidence
to collapse the relation.

\subsection{Isolation and collapse}

\begin{definition}[Isolation certificate]
At level \(k\), an isolation certificate for a candidate \(c\) consists of a
positive rational \(\delta\) and a proof that:
\begin{enumerate}[label=(\roman*)]
\item the region \(R(k,c)\) contains an intended solution;
\item every distinct candidate region is at distance at least \(\delta\);
\item future approximation error is less than \(\delta/3\).
\end{enumerate}
\end{definition}

\begin{theorem}[Relational collapse under isolation]
Given a finite relational presentation and an isolation certificate, bounded
search selects the unique compatible branch at every sufficiently fine
level.  The selected observations form a deterministic regular name.
\end{theorem}

\begin{proof}
At a fine level, compare each candidate region with the isolating region
using rational distance lower and upper bounds.  The true branch has distance
at most the accumulated approximation error, less than \(\delta/3\).
Every other branch has distance at least
\(\delta-\delta/3=2\delta/3\).  The rational tests are separated, so bounded
search returns exactly one branch.  Compatibility and the decreasing error
budget show that the selected branch at level \(k+1\) refines the branch at
level \(k\).  The centers therefore satisfy the regular-name estimate.
All alternatives are finite and all comparisons are exact cross-multiplied
rational inequalities.
\end{proof}

This theorem is the standard repair for exact Puiseux coefficients, rotation
axes, eigenvectors, and local inverse branches: first prove a uniform
relational theorem, then obtain a single exact branch only under explicit
isolation data.
